\documentclass[11pt,a4paper]{article}
\usepackage[margin=25mm,headheight=14pt]{geometry}
\usepackage{fontspec}
\setmainfont{lmroman10-regular.otf}[BoldFont=lmroman10-bold.otf,ItalicFont=lmroman10-italic.otf,BoldItalicFont=lmroman10-bolditalic.otf]
\setsansfont{lmsans10-regular.otf}
\setmonofont{lmmono10-regular.otf}[Scale=MatchLowercase]
\newfontfamily\authorsymbol{Noto Sans Symbols}[Scale=MatchUppercase]
\usepackage{amsmath,amssymb,amsthm,mathtools,booktabs}
\usepackage{microtype}
\usepackage[hidelinks,unicode]{hyperref}
\usepackage{fancyhdr}
\pagestyle{fancy}\fancyhf{}\fancyhead[L]{\small\itshape Exact cyclic-shift energies}\fancyhead[R]{\small\thepage}\renewcommand{\headrulewidth}{0pt}
\setlength{\parindent}{1.2em}\setlength{\parskip}{0pt}\setlength{\emergencystretch}{2em}
\newtheorem{theorem}{Theorem}\newtheorem{proposition}[theorem]{Proposition}\newtheorem{lemma}[theorem]{Lemma}\newtheorem{corollary}[theorem]{Corollary}\theoremstyle{definition}\newtheorem{definition}[theorem]{Definition}\newtheorem{example}[theorem]{Example}
\DeclareMathOperator{\rank}{rank}\DeclareMathOperator{\spanop}{span}\DeclareMathOperator{\diag}{diag}
\newcommand{\R}{\mathbb R}\newcommand{\C}{\mathbb C}\newcommand{\Z}{\mathbb Z}\newcommand{\N}{\mathbb N}\newcommand{\ip}[2]{\langle #1,#2\rangle}\newcommand{\norm}[1]{\lVert #1\rVert}\newcommand{\one}{\mathbf 1}
\hypersetup{pdftitle={Exact Cyclic-Shift Energies: Minimal State and Integral Bounds},pdfauthor={lmtlssss🜐},pdfsubject={Mathematics: harmonic analysis and discrete optimization},pdfkeywords={cyclic shifts, quadratic forms, circulant matrices, exact state, integer occupancy}}
\title{\LARGE Exact Cyclic-Shift Energies:\\Minimal State and Integral Bounds}
\author{lmtlssss{\authorsymbol 🜐}}
\date{2 November 2026\\[3pt]\small\href{https://doi.org/10.5281/zenodo.23897127}{10.5281/zenodo.23897127}}
\begin{document}\maketitle\thispagestyle{plain}

\begin{abstract}
We give an exact representation of cyclic-shift quadratic costs that separates
linear response state from pairwise interactions. For real vector-valued profiles
on $\Z/P\Z$ and a positive semidefinite shift-invariant quadratic form, the minimum
dimension of a linear summary preserving the cost of every real superposition is
$\sum_{k=0}^{P-1}\rank G_k$, where $G_k$ is the frequency-wise Gram matrix.
For a common profile with fixed real weights, one correlation kernel and one
retained potential determine every single-coordinate cost and admit an $O(P)$
exact update. For unweighted copies, integer occupancy gives a spectral lower
bound with the sharp residue term $r(P-r)/P$, where $r=N\bmod P$.
A spectrum-flat family attains the bound exactly at the balanced occupancies.
The paper includes proofs, integer and rational verification, and a
source-derived figure.
\end{abstract}
\noindent\textit{Keywords.} Cyclic shifts; quadratic forms; sufficient linear
state; circulant matrices; discrete optimization; integer occupancy.

\section{Cyclic profiles and the quadratic cost}
Let $P\ge2$ and let $\mathcal V=(\R^m)^P$. Indices are taken modulo $P$.
For $d\in\Z/P\Z$, define $(T_dx)(t)=x(t-d)$.
Let $A:\mathcal V\to\mathcal V$ be symmetric positive semidefinite and commute
with every $T_d$. The associated seminorm is
$\norm{x}_A^2=x^\mathsf T A x$.
For fixed profiles $x_1,\ldots,x_N\in\mathcal V$ and shift choices
$d=(d_1,\ldots,d_N)$, put
\begin{equation}\label{eq:cost}
 E(d)=\left\|\sum_{j=1}^N T_{d_j}x_j\right\|_A^2.
\end{equation}
We use the unitary discrete Fourier transform
\[
 \widehat x(k)=P^{-1/2}\sum_{t=0}^{P-1}x(t)e^{-2\pi i kt/P}.
\]
After complexification, $A$ has a Hermitian positive semidefinite symbol
$A_k\in\C^{m\times m}$ at frequency $k$. Real coefficients give conjugate
symmetry between $k$ and $P-k$. Define
\begin{equation}\label{eq:gram}
 X_k=[\widehat x_1(k)\ \cdots\ \widehat x_N(k)],\qquad
 G_k=X_k^*A_kX_k.
\end{equation}
The Fourier basis is a standard diagonalization of cyclic translation;
see Gray~\cite{gray}. Our first result specifies the exact state dimension
for a precise linear-summary contract.

\section{Minimum linear state for all superpositions}
Allow arbitrary real coefficients $c_{j,d}$ and form
\[
 Lc=\sum_{j=1}^N\sum_{d=0}^{P-1}c_{j,d}T_dx_j,
 \qquad Q(c)=\norm{Lc}_A^2.
\]
A linear summary is a linear map $S:\R^{NP}\to\R^q$ for which a function
$f$ satisfies $Q(c)=f(Sc)$ for every $c\in\R^{NP}$. The decoder $f$ may be
arbitrary. This definition states both the retained input family and the
information required from the summary.

\begin{theorem}[Minimum linear dimension]\label{thm:rank}
The least possible $q$ is
\begin{equation}\label{eq:rank}
 \boxed{q_{\min}=\rank(L^\mathsf T A L)
       =\sum_{k=0}^{P-1}\rank G_k.}
\end{equation}
A linear summary of that dimension can be updated additively after any
change of the superposition coefficients.
\end{theorem}
\begin{proof}
Write $B=A^{1/2}L$. Then $Q(c)=\norm{Bc}_2^2$ and
$\rank(B)=\rank(B^\mathsf T B)$.
If $Sc=0$, the summary contract gives
$\norm{Bc}_2^2=f(0)=Q(0)=0$. Hence $\ker S\subseteq\ker B$ and
$\rank S\ge\rank B$. Conversely, choose orthonormal coordinates on
$\operatorname{im}B$. These coordinates of $Bc$ give a summary with
$\rank B$ real entries, and their squared Euclidean norm recovers $Q$.
The coordinates are linear, so coefficient changes give additive updates.

For the frequency formula, apply a Fourier transform to the cyclic
coefficient sequence $(c_{j,d})_{d=0}^{P-1}$ for each $j$. At frequency $k$,
$L$ becomes $\sqrt P X_k$, and $B$ becomes
$\sqrt P A_k^{1/2}X_k$. Thus the complexified rank is the sum of these
block ranks. Each block has rank $\rank(X_k^*A_kX_k)$.
Complexification preserves the rank of a real matrix. Equivalently,
conjugate frequency pairs contribute twice their complex rank as a real
subspace; frequencies $0$ and, for even $P$, $P/2$ contribute their real
ranks once. This proves~\eqref{eq:rank}.
\end{proof}

The contract includes all real superpositions of translated profiles.
For a finite set of shift assignments alone, the information model is
different. Theorem~\ref{thm:rank} gives the minimum for a linear encoder
under the stated all-superposition contract.

\begin{corollary}[Scalar common profile]\label{cor:scalar}
For $m=1$, a single common profile $p$ and scalar symbol $A_k\ge0$, the
minimum dimension is the number of frequencies satisfying
$A_k|\widehat p(k)|^2>0$. In particular, a positive symbol and a profile
with all Fourier coefficients nonzero require exactly $P$ real coordinates.
\end{corollary}
\begin{proof}
The one-profile Gram block is the scalar $A_k|\widehat p(k)|^2$.
Its rank is either zero or one. Apply Theorem~\ref{thm:rank}.
\end{proof}

\section{One retained potential for exact coordinate updates}
Suppose $x_j=a_jp$ for fixed real weights $a_j$ and a common profile
$p\in\mathcal V$. Define the real symmetric correlation kernel
\begin{equation}\label{eq:kernel}
 g(u)=\ip{p}{T_up}_A,
 \qquad g(-u)=g(u),
\end{equation}
and the retained potential
\begin{equation}\label{eq:potential}
 \Phi_d(s)=\sum_{j=1}^N a_jg(d_j-s).
\end{equation}
The notation $\ip{x}{y}_A=x^\mathsf T A y$ is used throughout.

\begin{theorem}[Coordinate-cost and update identities]\label{thm:update}
For every shift vector and every real weight vector,
\begin{align}
 E(d)&=\sum_i a_i\Phi_d(d_i),\label{eq:ephi}\\
 E(d_1,\ldots,s,\ldots,d_N)
 &=E_{-i}+a_i^2g(0)+2a_i\bigl[\Phi_d(s)-a_i g(d_i-s)\bigr],\label{eq:coord}
\end{align}
where $E_{-i}=\norm{\sum_{j\ne i}a_jT_{d_j}p}_A^2$.
If coordinate $i$ changes from $u$ to $v$, then
\begin{equation}\label{eq:phiupdate}
 \Phi_{d'}(s)=\Phi_d(s)+a_i\{g(v-s)-g(u-s)\}.
\end{equation}
After construction of $g$ and $\Phi$, evaluating all $P$ choices of one
coordinate and updating the retained potential cost $O(P)$ arithmetic
operations. The kernel and potential occupy $2P$ scalar entries, in
addition to $O(N)$ weights and shifts.
\end{theorem}
\begin{proof}
Translation invariance gives
$\ip{T_up}{T_vp}_A=g(v-u)$. Expanding the square proves
\eqref{eq:ephi}. Separating the $i$th term into its self contribution and
its inner product with the remaining sum gives~\eqref{eq:coord}.
Only one summand in~\eqref{eq:potential} changes, which gives
\eqref{eq:phiupdate}. Each displayed formula needs a constant number of
scalar operations at each cyclic position. The two length-$P$ arrays
and the weights and shifts give the stated storage.
\end{proof}

The preprocessing cost is explicit. For a scalar profile, first apply $A$ to $p$, charging that operation
explicitly; the $P$ direct correlations then require $O(P^2)$ arithmetic
operations. A Fourier implementation instead uses the known symbol and
its transform costs. Directly
forming~\eqref{eq:potential} costs $O(NP)$ operations. These are separate
from the $O(P)$ cost of an accepted move. A table storing all $N^2$
pair kernels contains $N^2P$ scalar entries; the common-profile
factorization replaces that table by one kernel.

A coordinate method can choose the smallest minimizing shift under a
fixed tie rule. Strict-decrease moves terminate because the shift space
is finite. The identities are exact objective evaluations; a fixed point
of coordinate search is a coordinatewise minimum. A global certificate
is supplied by the next section for a specified spectral family.

\section{Integer occupancy and a sharp spectral bound}
Take $a_j=1$ and define
$h_u=|\{j:d_j=u\}|$. Then $h\in\Z_{\ge0}^P$, $\one^\mathsf Th=N$, and
\[
 E=h^\mathsf TKh,\qquad K_{uv}=g(v-u).
\]
The real symmetric circulant matrix $K$ is positive semidefinite.
Let $\lambda_0$ be its eigenvalue on $\one$, and let
$\lambda_*\ge0$ be its smallest eigenvalue on $\one^\perp$.
Write $N=qP+r$, with $q\ge0$ and $0\le r<P$.

\begin{lemma}[The integer residue term]\label{lem:balance}
Every nonnegative integer occupancy of total mass $N$ satisfies
\begin{equation}\label{eq:integer}
 \left\|h-\frac NP\one\right\|_2^2
 \ge\frac{r(P-r)}P.
\end{equation}
Equality holds precisely when $r$ entries equal $q+1$ and the remaining
$P-r$ entries equal $q$.
\end{lemma}
\begin{proof}
If two entries differ by at least two, transferring one unit from the
larger to the smaller reduces $\sum h_u^2$ by
$2(h_{\max}-h_{\min}-1)>0$. Repeating this operation ends at entries that
differ by at most one. Their sum forces the stated multiplicities.
Their squared norm is $Pq^2+2qr+r$. Subtracting $N^2/P$ gives
$r(P-r)/P$. Conversely, any unbalanced vector admits the strict reduction.
\end{proof}

\begin{theorem}[Sharp integral spectral bound]\label{thm:occupancy}
For every shift assignment,
\begin{equation}\label{eq:lower}
 \boxed{E\ge\frac{\lambda_0N^2}{P}
              +\lambda_*\frac{r(P-r)}P.}
\end{equation}
If $K=\alpha I+\beta\one\one^\mathsf T$ with $\alpha>0$ and $\beta\ge0$,
the exact minimum is
\begin{equation}\label{eq:flat}
 E_{\min}=\alpha(Pq^2+2qr+r)+\beta N^2,
\end{equation}
and its occupancies are exactly the $\binom Pr$ balanced vectors.
For a general $K$ with $\lambda_*>0$, equality in~\eqref{eq:lower} holds
precisely when the occupancy is balanced and its centered vector belongs
to the $\lambda_*$ eigenspace (with the zero vector allowed).
\end{theorem}
\begin{proof}
Set $v=h-(N/P)\one$. Since $v\perp\one$ and $K\one=\lambda_0\one$,
\[
 h^\mathsf TKh=\lambda_0N^2/P+v^\mathsf TKv
 \ge\lambda_0N^2/P+\lambda_*\norm v_2^2.
\]
Lemma~\ref{lem:balance} proves~\eqref{eq:lower}. In the flat family,
$\lambda_0=\alpha+\beta P$ and $K$ acts by $\alpha$ on $\one^\perp$;
all balanced occupancies therefore attain~\eqref{eq:flat}.
For positive $\lambda_*$, equality requires equality in both inequalities,
which gives the final characterization.
\end{proof}

If $P$ divides $N$, uniform occupancy removes every nonconstant Fourier
coefficient and attains the constant-mode lower bound for every such
positive semidefinite $K$. When $\lambda_*>0$, this occupancy is unique.
The residue term describes exactly the integrality obstruction to that
uniform assignment when $P$ does not divide $N$.

\section{Exact examples and verification}
Take $P=16$, $N=64$, $A=I$, and
$p=(2,1,0,\ldots,0)$. Its circulant correlation symbol is
\[
 \lambda_k=|2+e^{-2\pi i k/16}|^2
 =5+4\cos(2\pi k/16).
\]
Every eigenvalue is positive; the constant eigenvalue is nine and the
smallest eigenvalue is one. Corollary~\ref{cor:scalar} therefore gives a
minimum all-superposition linear summary of sixteen real coordinates.

Put four copies at every shift. Each output position receives four
contributions of weight two and four of weight one. Thus
\[
 h=(4,\ldots,4),\qquad
 \sum_jT_{d_j}p=(12,\ldots,12),\qquad E=16\cdot12^2=2304.
\]
Theorem~\ref{thm:occupancy} gives the same lower bound,
$\lambda_0N^2/P=9\cdot64^2/16=2304$, so this assignment is globally
optimal. The occupancy is unique, while permutations of labelled copies
produce multiple shift vectors with that occupancy.

If all copies are aligned, the only nonzero output values are $128$ and
$64$, giving $E=128^2+64^2=20480$. These are two exact evaluations of the
same quadratic objective on the same profile family.

The standard pair-kernel table has $64^2\cdot16=65536$ scalar entries.
The common-profile kernel and retained potential use $2\cdot16=32$.
Weights and shifts occupy their separate $O(N)$ arrays, and construction
cost is charged as in Theorem~\ref{thm:update}. The sixteen-coordinate
minimum and the thirty-two-entry implementation refer respectively to
the linear-summary contract and the named kernel-plus-potential arrays.

The companion verifier uses integer and rational arithmetic. It checks
6822 coordinate costs against direct profile summation and filtering,
1200 retained-potential updates against reconstruction, and 19437 integer
occupancies against~\eqref{eq:flat}. Thirty-two independently generated
two-component, four-position examples check~\eqref{eq:rank} using exact
Gaussian-integer matrix ranks. The figure is generated from the
$P=16$, $N=64$ exact occupancy certificate.

\section{Relation to established methods}
Fourier diagonalization of circulant forms is classical~\cite{gray}.
Quadratic optimization over unit-modulus variables and coordinate
methods have a substantial literature; Ragi, Chong and Mittelmann
provide approximation methods and guarantees for a related unimodular
problem~\cite{ragi}. Here a coordinate is a cyclic shift, which ties its
phase across all frequencies. The results specify an exact linear-state
contract, a complete retained-coordinate identity with charged storage,
and a sharp integer-occupancy certificate. Their proofs are independent
of an approximation or a claim of general global coordinate convergence.

\begin{thebibliography}{9}
\bibitem{gray} R. M. Gray, \emph{Toeplitz and Circulant Matrices: A Review},
Foundations and Trends in Communications and Information Theory 2(3)
(2006), 155--239. \href{https://doi.org/10.1561/0100000006}{doi:10.1561/0100000006}.
\bibitem{ragi} S. Ragi, E. K. P. Chong and H. D. Mittelmann,
Polynomial-time methods to solve unimodular quadratic programs with performance
guarantees, IEEE Transactions on Aerospace and Electronic Systems 55(5)
(2019), 2118--2127. \href{https://doi.org/10.1109/TAES.2018.2881366}{doi:10.1109/TAES.2018.2881366}.
\end{thebibliography}
\end{document}
